Dependability in space-grade electronics has long been a demanding discipline: radiation-induced single-event upsets can silently corrupt registers, memories, and --- in \acp{FPGA} --- the configuration fabric itself, altering the circuit behaviour rather than just its state. \ac{RISC-V}'s openness and the availability of synthesisable \ac{SoC} ecosystems have made it an attractive foundation for space-oriented platforms, but the freedom to add fault-tolerance mechanisms brings a new challenge: choosing the right ones requires comparative data that is rarely available, because no standard platform exists to integrate, stress, and measure them under equal conditions. This thesis introduces \ac{Fatori-V}, a configurable \ac{RISC-V} \ac{SoC} platform that closes that gap by combining hardware fault-tolerance mechanisms, a standalone fault-injection tool, and an automated experiment workflow into a single reproducible system.